% 第二章：进阶方法（进阶篇）
% 示范带算法伪代码与表格的章节写法。
\chapter{进阶方法}

\section{改进的优化算法}
% 引出比第一章更一般的算法框架
在基础梯度下降之上，引入动量项可以加速收敛。
更新规则如式~\eqref{eq:momentum}：
\begin{equation}
  \vect{v}_{t+1} = \beta \vect{v}_t - \eta \nabla \mathcal{L}(\vect{\theta}_t),
  \qquad
  \vect{\theta}_{t+1} = \vect{\theta}_t + \vect{v}_{t+1}.
  \label{eq:momentum}
\end{equation}
其中 $\beta \in [0, 1)$ 为动量系数，$\eta$ 为学习率。

\section{超参数的影响}
% 用表格对比不同设置下的表现
表~\ref{tab:hyper} 给出了不同动量系数下的迭代次数对比。

\begin{table}[!htbp]
  \centering
  \caption{不同动量系数下的收敛迭代次数}
  \label{tab:hyper}
  \begin{tabular}{lcccc}
    \toprule
    动量系数 $\beta$ & 0 & 0.5 & 0.9 & 0.99 \\
    \midrule
    收敛迭代次数 & 812 & 356 & 154 & 203 \\
    \bottomrule
  \end{tabular}
\end{table}

\section{本章小结}
动量方法通过累积历史梯度方向改善收敛速度，
适合目标函数存在狭长山谷形等值面的场景；
实际使用中建议从 $\beta = 0.9$ 开始尝试。
